\section{行政主体}
\begin{enumerate}
  \item 解决控权对象问题。只看在具体的管理关系中。
  \item 市。省市（省、自治区、直辖市）、地级市。县级市 = 县（\al{昆山市}）。
  \item 内设机构。部司处。部 \yjt 对外。是公法人，吸收了内设机构和公务员的人格。司处 \yjt 内设机构。
  \item 行政机关分类。
        \begin{enumerate}
          \item 一般行政机关。政府（综合性）。国省市县乡 5 级建制。乡无对外独立部门。
          \item 专门行政机关。又称政府的（职能）部门，政府下属只负责某一个领域的。
                \begin{enumerate}
                  \item 议事协调机构。牵桥搭线的。\al{防震减灾、禁毒委员会}。
                  \item 临时机构。\al{疫情指挥部、联合调查组}。\sjt 这些都不属于行政主体。
                \end{enumerate}
          \item 正式行政机关。政府 + 部门。
          \item 派出机关。不怎么考。
                \begin{enumerate}
                  \item 街道办\ydy 乡政府。区政府派出。
                  \item 行政公署 \ydy 市。省政府派出。
                  \item 区公所。县政府派出，已经绝种了。
                  \item vs 居委会。是基层自治组织。
                \end{enumerate}
          \item 派出机构。部门派出的内设机构。原则上不能做行政主体，除非有授权。
                \begin{enumerate}
                  \item 派出所能做 5 警。500 以下罚款、警告。
                  \item 税务所、派出法庭。
                \end{enumerate}
          \item 开展行政活动的其他组织。
                \begin{enumerate}
                  \item 后天公法人 = 法律、法规、规章授权无行政权能组织。其他规范性文件无授权资格。
                  \item vs 委托。是委托关系，不是公法人，
                \end{enumerate}
                \begin{kk}{委托 vs 授权}
                  \begin{enumerate}
                    \item 这事我能委托还是授权给你我说了不算，得看法律、法规、规章咋说的。
                    \item 不同的权力有资格授予这个权力的文件不同（谁给你这个权力去授权的）。
                          \begin{enumerate}
                            \item 强制措施权 \yjt 只能法律、行政法规。
                            \item 普通权力（有单行立法） \yjt 法律、法规。\al{处罚、许可}。
                            \item 一般权力（没有单行立法） \yjt 法律、法规、规章。\al{给付、征收、裁决}。
                          \end{enumerate}
                  \end{enumerate}
                \end{kk}
        \end{enumerate}
  \item 行政主体的判定。权、名、责。蒙猜办法，“名字不正经”的一般都不是。
  \item 机构设置。\al{机构名称、职能 \dui}。题目会告诉你是啥类型。
        \begin{enumerate}
          \item 设立。一高一低留中间。
                \begin{enumerate}
                  \item 一高（组成部门 \yjt 部委行署 \yjt 全人大/常。\al{xx 委员会 \dui。xx \xhx{监督管理}委员会 \cuo}。
                  \item 一低（处级内设机构） \yjt 部委自己。\al{司法考试改革研究处 \yjt 司法部自己}。
                  \item 留中间 \yjt 统统交给国务院。\al{税务总局为国务院直属机构、禁毒委员会是国务院议事协调机构}。
                  \item 中有决定、批准（司级内设机构）。孙子辈国务院决定，儿子辈批准。\al{高教司}。
                  \item 国务院编制管理机关 \bdy 国务院！！
                \end{enumerate}
          \item 职能。组基直专，办事协助，部管特定。
                \begin{enumerate}
                  \item 组成部门履行基本职能。国计民生。看到组成部门要看全，可能是部管局。
                  \item 直属机构履行专门职能。\al{税务总局，旅游总局}。
                  \item 办事机构（内设机构）协助总理办公。
                  \item 部管局主管特定业务，行使行政管理职能。\al{国家海洋局为国务院组成部门管理的国家局}。
                \end{enumerate}
          \item 地方。上下级编管是指导关系，不是领导关系。
                \begin{enumerate}
                  \item 变闺名，本级提，上级批。变更规格、名称。本级政府（编管）提方案 \yjt 上级政府（编管）批准。
                  \item 设、撤、并，还要同人备。设立、撤销、合并（+ 同级人大备案）。
                  \item 内设机，本级定。内设机构 \yjt 本级政府（编管）决定。
                  \item 异议商，本级调，本级定。异议 \yjt 协商 \yjt 本级政府编管协调 \yjt 本级政府决定。
                  \item 协调不用上来批。议事协调机构无需报上一级政府批准。
                \end{enumerate}
        \end{enumerate}
  \item 编制管理。针对人。\al{机构人员的数量定额、领导职数 \dui}。后期死背就行了。
        \begin{enumerate}
          \item 基本原则。中央统一领导、地方分级管理。
                \begin{enumerate}
                  \item 应当确定、不得干预、要求对口。
                  \item 设立机构时应当单独确定人员、编制。
                  \item 议事协调机构。
                        \begin{enumerate}
                          \item 没有编制 = 不单独确定编制。牵桥搭线的，人员都来自各部门。
                          \item 议定事项经国务院同意，有关的行政机构按各自的职责负责。
                          \item \al{可以规定行政措施、可以设立司、处两级内设机构 \cuo}。
                        \end{enumerate}
                \end{enumerate}
          \item 国务院。可自己决定。两批、两跨、一审。
                \begin{enumerate}
                  \item 批准各级政府的总额、批准政法、公安、国安、监狱等特别部门的专项名额。
                  \item 国务院编管拍板跨部门、跨层次调整编制。
                  \item 形式审核地方编制。
                  \item 区分批准、调整、“跨”调整。\al{甲\xhx{市政府}可以在\xhx{总额内}调整\xhx{市政府}有关部门的行政编制 \dui}。
                \end{enumerate}
          \item 政府。
                \begin{enumerate}
                  \item 省政府。发布地方事业单位编制。
                  \item 地方政府。\yi 编制管理是地方事务，发布 \yjt 上一级政府。\er 权限争议 \yjt 上级。
                \end{enumerate}
        \end{enumerate}
  \item 评估应当定期进行。
\end{enumerate}
